% !PRINTED 199
\weberpage{199}
%<<CONT
\gesperrt{corresponds to a multiplication. The transforming substitution $L$ is thereby
itself determined up to a multiplication.}

\begin{center}
\S.\,52.

\gesperrt{The different kinds of possible groups.}
\end{center}

Let $a$ be one of the poles of the group $G$, and let us
assume that there are $\nu - 1$ and no more elements in $G$,
$\Theta_1, \Theta_2, \ldots, \Theta_{\nu-1}$, such that
\weq{1}{a = \Theta_1(a) = \Theta_2(a) \cdots = \Theta_{\nu-1}(a)}
holds. Such a pole shall be called a $\nu$-fold pole. It is
now first of all clear that the elements
\weq{2}{1, \Theta_1, \Theta_2, \ldots, \Theta_{\nu-1}}
form a group by themselves, and indeed a divisor of $G$;
for from
\[
a = \Theta_1(a), \quad a = \Theta_2(a)
\]
it follows, if on the right-hand side of the second equation
one replaces $a$ by $\Theta_1(a)$, which is equal to it:
\[
a = \Theta_2 \Theta_1(a).
\]
\gesperrt{Hence $\nu$ must be a divisor of $n$:}
\weq{3}{n = \nu\mu.}
We denote the group (2) by $Q$.

It can be proved that this group must be cyclic, that is,
consists of the powers of one of its elements.

For if by transformation according to Theorem 2.
we send $a$ to infinity, the substitutions (2) take the
form:
\[
x, \; \varepsilon_1 x + c_1, \; \varepsilon_2 x + c_2, \ldots, \varepsilon_{\nu-1} x + c_{\nu-1},
\]
and the composition of two among them gives:
\[
\begin{aligned}
\Theta_1 \Theta_2 &= \varepsilon_1 \varepsilon_2 x + (c_2 \varepsilon_1 + c_1) \\
\Theta^\lambda &= \varepsilon^\lambda x + c(\varepsilon^{\lambda-1} + \varepsilon^{\lambda-2} + \cdots + 1) \\
\Theta^{-1} &= \varepsilon^{-1} x - c \varepsilon^{-1}.
\end{aligned}
\]
From this one concludes that all $\varepsilon_1, \varepsilon_2, \ldots, \varepsilon_{\nu-1}$ must be roots
of unity of degree $\nu$. Moreover, no two of the $\varepsilon$ can be
equal to one another. For if, e.g., $\varepsilon_1 = \varepsilon_2$, then we should have
\[
\Theta_1 \Theta_2^{-1} = x + (c_1 - c_2).
\]
